\section{Results}
\label{sec:results}

% ════════════════════════════════════════════════════════════════════════════
\subsection{In evolution, duplicated bodies end ahead of tree bodies}
\label{sec:evolution}

\begin{figure}[t]
  \centering
  \includegraphics[width=\linewidth]{evolution_genomes.png}
  \caption{Evolution with the \bb\ controller, 20 seeds per genome. (a) Best score so far in each run, mean $\pm$ SEM over seeds. (b) Paired difference to the tree run of the same seed, with a 95\% bootstrap CI.}
  \label{fig:evolution}
\end{figure}

We first evolve bodies and controllers together (\cref{fig:evolution}, \cref{tab:evolution}). With the \bb\ controller, tree runs end at \SI{3.90}{m}, symmetric runs at \SI{4.35}{m} and tied runs at \SI{4.44}{m}. The tied lead over tree is clear ($+0.54$~m, 95\% CI $[+0.26, +0.82]$, $p = 0.004$). The symmetric lead is about as large but much more variable across seeds, and its interval includes zero ($+0.45$~m, $[-0.12, +1.04]$, $p = 0.25$). Symmetric and tied runs do not differ. Symmetric runs start ahead, by $+0.40$~m in the initial population, and end with a similar lead after a dip mid-run, while tied runs start level with tree and pull ahead over the generations.

Evolved tied bodies are also larger than the others, and 8 of the 20 tied champions exceed the module limit. Size does not explain the \bb\ result: with those eight seeds removed, tied runs still lead tree runs by $+0.48$~m ($[+0.15, +0.80]$, $p = 0.03$, 12 seeds).

\begin{table}[t]
  \centering
  \small
  \caption{Evolution, 20 seeds per arm. Left: mean best score (m). Right: paired gaps (m) with bootstrap 95\% CI, Wilcoxon $p$ and the number of seeds in which the first genome is ahead.}
  \label{tab:evolution}
  \resizebox{\linewidth}{!}{\begin{tabular}{lccc@{\hskip 1.2em}lcc}
    \toprule
    controller & \tree & \sym & \tied & comparison & gap [95\% CI] & $p$ (wins) \\
    \midrule
    \multirow{3}{*}{\bb}     & \multirow{3}{*}{3.90} & \multirow{3}{*}{4.35} & \multirow{3}{*}{4.44} & symmetric $-$ tree & $+0.45$ $[-0.12, +1.04]$ & 0.25 (11) \\
                             &  &  &  & tied $-$ tree      & $+0.54$ $[+0.26, +0.82]$ & 0.004 (15) \\
                             &  &  &  & symmetric $-$ tied & $-0.09$ $[-0.59, +0.41]$ & 0.73 (9) \\
    \midrule
    \multirow{3}{*}{\bbdrift} & \multirow{3}{*}{3.31} & \multirow{3}{*}{3.66} & \multirow{3}{*}{3.45} & symmetric $-$ tree & $+0.35$ $[-0.05, +0.79]$ & 0.15 (11) \\
                             &  &  &  & tied $-$ tree      & $+0.14$ $[-0.15, +0.40]$ & 0.23 (13) \\
                             &  &  &  & symmetric $-$ tied & $+0.21$ $[-0.20, +0.62]$ & 0.28 (13) \\
    \midrule
    \multirow{3}{*}{\bbfree} & \multirow{3}{*}{1.45} & \multirow{3}{*}{1.58} & \multirow{3}{*}{1.53} & symmetric $-$ tree & $+0.13$ $[+0.01, +0.24]$ & 0.053 (14) \\
                             &  &  &  & tied $-$ tree      & $+0.08$ $[-0.02, +0.21]$ & 0.19 (13) \\
                             &  &  &  & symmetric $-$ tied & $+0.05$ $[-0.09, +0.16]$ & 0.39 (12) \\
    \bottomrule
  \end{tabular}}
\end{table}

% ════════════════════════════════════════════════════════════════════════════
\subsection{On fixed bodies, the symmetric lead builds during learning}
\label{sec:fixed}

\begin{figure}[t]
  \centering
  \includegraphics[width=0.92\linewidth]{fixed_body_curves.png}
  \caption{Learning on fixed random bodies, 80 bodies per set and three training seeds per body. Top: best score so far, mean $\pm$ SEM over bodies. Middle: difference to the tree mean at the same number of evaluations (band: bootstrap 95\% CI of symmetric $-$ tree). Bottom: mean of each set divided by the tree mean, from 60 evaluations. All sets have the hinge counts of the symmetric bodies.}
  \label{fig:curves}
\end{figure}

Evolution mixes which bodies are made with how well their controllers are trained. To separate the two, we train controllers on fixed random bodies with equal hinge counts (\cref{fig:curves}, \cref{tab:fixed}). After training, symmetric bodies walk \SI{0.37}{m} further than tree bodies under \bb\ ($+25\%$; 95\% CI $[+0.07, +0.67]$) and \SI{0.28}{m} further under \dann\ ($+18\%$; $[-0.05, +0.60]$). The effect is modest. Under \bb\ the permutation test on the means gives $p = 0.016$ and the rank test $p = 0.07$; under \dann\ neither reaches 0.05. A randomly chosen symmetric body beats a randomly chosen tree body 58\% of the time under \bb\ and 55\% under \dann.

The lead is not there at the start. After the first 20 evaluations, when nothing has been learned, the two genomes do not differ (\bb\ $+0.03$~m, $[-0.04, +0.09]$, $p = 0.44$; \dann\ $+0.02$~m, $[-0.02, +0.07]$, $p = 0.31$). It appears during learning: from 20 evaluations to the end of training, symmetric bodies gain more than tree bodies (\bb\ $+0.34$~m, $[+0.08, +0.60]$, $p = 0.013$; \dann\ $+0.26$~m, $[-0.04, +0.56]$, $p = 0.11$). The absolute gap keeps growing to the end of training, while the ratio to tree peaks early and then falls slightly, from 1.30 at 200 evaluations to 1.25 under \bb\ and from 1.25 to 1.18 under \dann.

The two \bb\ variants show that the lead needs a shared clock. \Bbdrift, whose hinge frequencies stay close to a shared one, keeps the lead ($+0.25$~m, $[+0.02, +0.48]$). \Bbfree, in which every hinge has its own frequency, loses it ($-0.02$~m, $[-0.12, +0.08]$), but only because every body does poorly: scores fall to \SI{0.63}{m} on fixed bodies and to about \SI{1.5}{m} in evolution (\cref{tab:fixed,tab:evolution}, \cref{fig:curves-variants} in the appendix).

\begin{table}[t]
  \centering
  \small
  \caption{Symmetric and tree on 80 + 80 fixed random bodies, after 1{,}600 evaluations (\dann: 5{,}000). Seeds are averaged per body. Scores and gaps in metres; gap with bootstrap 95\% CI over bodies, permutation $p$ on the mean, Mann--Whitney $p$, and the ratio of means.}
  \label{tab:fixed}
  \begin{tabular}{lcccccc}
    \toprule
    controller & \tree & \sym & gap [95\% CI] & perm.\ $p$ & M--W $p$ & ratio \\
    \midrule
    \bb      & 1.44 & 1.81 & $+0.37$ $[+0.07, +0.67]$ & 0.016 & 0.070 & 1.25 \\
    \bbdrift & 1.22 & 1.47 & $+0.25$ $[+0.02, +0.48]$ & 0.037 & 0.13 & 1.20 \\
    \bbfree  & 0.65 & 0.63 & $-0.02$ $[-0.12, +0.08]$ & 0.71 & 0.38 & 0.97 \\
    \dann    & 1.58 & 1.85 & $+0.28$ $[-0.05, +0.60]$ & 0.096 & 0.28 & 1.18 \\
    \bottomrule
  \end{tabular}
\end{table}

% ════════════════════════════════════════════════════════════════════════════
\subsection{Bodies without mirroring match symmetric bodies}
\label{sec:bodies}

\begin{figure}[t]
  \centering
  \includegraphics[width=0.9\linewidth]{shape_measures.png}
  \caption{Resting shape of the five body sets (80 bodies each, after the \SI{3}{s} settle). Boxes: quartiles; dots: bodies.}
  \label{fig:shape}
\end{figure}

At equal hinge count, random symmetric bodies are shaped differently from tree bodies (\cref{fig:shape}). They are flatter (height 0.14 vs \SI{0.23}{m}), spread their floor contacts further front to back (0.15 vs \SI{0.12}{m}), and have more yaw-axis leg hinges (1.8 vs 1.1) and fewer lateral-axis hinges (0.6 vs 1.7). Does this shape carry the lead, or does the mirroring itself? We test this with three sets of bodies that are not mirrored (\cref{fig:body-sets}, \cref{tab:sets}).

Tree bodies chosen to match the symmetric bodies' resting shape do as well as symmetric bodies. They beat ordinary tree bodies by $+0.33$~m under \bb\ ($p = 0.010$) and $+0.30$~m under \dann\ ($p = 0.048$), and the difference from symmetric bodies is close to zero ($+0.03$ and $-0.02$~m). Regression agrees: adjusting the symmetric $-$ tree gap for the six shape measures reduces it from $+0.37$ to $+0.08$~m under \bb\ and from $+0.28$ to $-0.03$~m under \dann.

Tied bodies, which duplicate modules without mirroring them, also do as well as symmetric bodies, and they beat tree bodies by $+0.29$~m under \bb\ ($p = 0.034$) and $+0.37$~m under \dann\ ($p = 0.016$). They reach this without the symmetric shape: their height falls between tree and symmetric bodies (\SI{0.17}{m}), and their lateral-axis and yaw-axis hinge counts stay close to those of tree bodies. Having more legs is not the explanation either. Tied bodies have more leg hinges than tree bodies (3.5 vs 2.9), but tree bodies matched to them on leg hinges gain almost nothing ($+0.09$ and $+0.03$~m). What tied bodies gain from duplication remains open.

These comparisons show that mirroring is not required for the lead. They do not show that mirroring adds nothing: the intervals on symmetric $-$ shape-matched and symmetric $-$ tied reach about $\pm 0.3$~m, as large as the lead itself.

\begin{figure}[t]
  \centering
  \includegraphics[width=0.75\linewidth]{body_sets.png}
  \caption{Final score per body set (mean $\pm$ SEM over 80 bodies, seeds averaged per body).}
  \label{fig:body-sets}
\end{figure}

\begin{table}[t]
  \centering
  \small
  \caption{Final score per body set (m), and each set's gap to tree and to symmetric bodies (m), with bootstrap 95\% CI over bodies and permutation $p$ on the mean.}
  \label{tab:sets}
  \resizebox{\linewidth}{!}{\begin{tabular}{llcll}
    \toprule
    controller & body set & mean & gap to tree [95\% CI], $p$ & symmetric $-$ set [95\% CI], $p$ \\
    \midrule
    \multirow{5}{*}{\bb} & tree & 1.44 & --- & $+0.37$ $[+0.07, +0.67]$, 0.016 \\
    & tree, leg-matched & 1.53 & $+0.09$ $[-0.18, +0.34]$, 0.50 & --- \\
    & tree, shape-matched & 1.78 & $+0.33$ $[+0.08, +0.59]$, 0.010 & $+0.03$ $[-0.27, +0.34]$, 0.83 \\
    & tied & 1.73 & $+0.29$ $[+0.03, +0.54]$, 0.034 & $+0.08$ $[-0.23, +0.39]$, 0.61 \\
    & symmetric & 1.81 & $+0.37$ $[+0.07, +0.67]$, 0.016 & --- \\
    \midrule
    \multirow{5}{*}{\dann} & tree & 1.58 & --- & $+0.28$ $[-0.05, +0.60]$, 0.096 \\
    & tree, leg-matched & 1.61 & $+0.03$ $[-0.24, +0.31]$, 0.83 & --- \\
    & tree, shape-matched & 1.87 & $+0.30$ $[+0.02, +0.59]$, 0.048 & $-0.02$ $[-0.36, +0.31]$, 0.91 \\
    & tied & 1.95 & $+0.37$ $[+0.07, +0.67]$, 0.016 & $-0.09$ $[-0.44, +0.25]$, 0.60 \\
    & symmetric & 1.85 & $+0.28$ $[-0.05, +0.60]$, 0.096 & --- \\
    \bottomrule
  \end{tabular}}
\end{table}

% ════════════════════════════════════════════════════════════════════════════
\subsection{Symmetric gaits lock, but do not depend on their left--right timing}
\label{sec:gaits}

\begin{figure}[t]
  \centering
  \includegraphics[width=\linewidth]{locked.png}
  \caption{Phase locking and score. (a) Final score of trainings on fixed random bodies (mean $\pm$ SEM over trainings): tree bodies, and symmetric bodies split by whether the trained gait is locked. (b) Share of evolved champions that are locked, by controller and genome.}
  \label{fig:locked}
\end{figure}

Symmetric bodies also move differently. Their left and right limbs often keep a steady lag from cycle to cycle, which we call phase locking (\cref{sec:measures}). In evolution with \bb, 15 of 20 symmetric champions are locked, against 2 of 20 tree champions and 3 of 20 tied champions (\cref{fig:locked}b), so tied bodies reach symmetric scores without locking. On fixed bodies, 43\% of symmetric trainings under \bb\ and 50\% under \dann\ end locked, against 2\% of tree trainings. Locked gaits are the good ones (\cref{fig:locked}a): locked symmetric trainings score \SI{2.14}{m} under \bb, unlocked ones \SI{1.56}{m} and tree trainings \SI{1.44}{m}, and under \dann\ the same three groups score 2.12, 1.59 and \SI{1.58}{m}. On symmetric bodies that lock in some seeds but not others, the locked seeds score higher under \bb\ ($+0.22$~m, 28 bodies, $p = 0.003$), but not clearly under \dann\ ($+0.09$~m, 32 bodies, $p = 0.30$). Locking also goes with shape: bodies that lock in all three seeds under \bb\ are flatter than bodies that never lock (height 0.10 vs \SI{0.16}{m}, exploratory).

\begin{figure}[t]
  \centering
  \includegraphics[width=\linewidth]{phase_shift.png}
  \caption{Phase-shift replay on fixed random bodies. Score change as a share of each controller's unperturbed score, when the left-side hinges are phase-shifted (top) or detuned (bottom). Solid: left side; dashed: the same number of randomly chosen lateral hinges (mean of three draws). Mean $\pm$ SEM over controllers with an unperturbed score of at least \SI{0.05}{m}.}
  \label{fig:phase-shift}
\end{figure}

Does the lock make these gaits good? We test this by perturbing the timing of trained gaits (\cref{sec:phase-shift}). We delay or detune the left side of the body, and compare with applying the same perturbation to the same number of randomly chosen hinges. If a gait depends on its left--right timing in particular, perturbing the left side should cost more.

It does not (\cref{fig:phase-shift}). Every gait depends heavily on its timing: a shift of a twentieth of a cycle costs a fifth to a third of the score, and half a cycle costs about 90\%. But for locked symmetric gaits, perturbing the left side costs no more than perturbing random hinges. Under \bb\ the side $-$ random difference is positive for every perturbation, from $+0.05$ to $+0.24$~m (33 bodies, all $p \geq 0.16$), so if anything the left-side perturbation costs less. Under \dann\ it ranges from $-0.18$ to $+0.03$~m (36 bodies, all $p \geq 0.08$). Detuning by 15\% breaks most of the lock (locked share 0.88 to 0.21 under \bb, 0.90 to 0.22 under \dann), yet it costs no more than the random control, which breaks less of it (0.34 and 0.35): 75\% vs 81\% of the score under \bb\ and 66\% vs 72\% under \dann. Locked gaits are not more fragile than tree gaits either; a shift of 0.05 cycles costs them about 20\% of their score, against 35\% for tree gaits under \bb\ and 34\% under \dann.

Locking therefore goes with good gaits, but trained gaits do not depend on it. This test concerns finished gaits. It does not rule out that the lock helps a controller find a good gait during learning.
